\section{Worked Example: Using the Companion's Debug and Profiling Tools}
\index{debug tools}\index{profiling}\index{event logger}\index{TSX profiling}\index{race checker}
\label{sec:debug-profiling-tools}

The preceding sections built tests, fuzzers, schedulers, and model
checkers.  This section introduces three toolchains from the companion
repository that support a different phase of development: diagnosing why
a backend fails a test, understanding where it spends its cycles, and
catching missing TM annotations before they become data races.

\subsection{The STM Bug Tool}
\index{STM bug tool}\index{event log parser}\index{invariant checker}
\label{sec:stm-bug-tool}

When an invariant test fails, the raw symptom---``Money destroyed by
32''---is rarely enough.  The implementation must expose what happened
inside the failing transactions: which addresses were read and written,
in what order, and where the protocol diverged from correctness.  The
companion's \texttt{tools/stm\_bug\_tool/} provides two stages of
analysis.

\paragraph{Event parsing.}
The file \texttt{event\_parser.py}~(\rep{tools/stm\_bug\_tool/event\_parser.py})
parses the output of a binary compiled with \texttt{-DTM\_EVENT\_LOGGER}.
The logger emits one line per TM operation in a ring buffer that the
tool flattens into structured records.  The core pattern is:

\begin{lstlisting}[language=Python, caption={Event regex from event\_parser.py.}]
EVENT_RE = re.compile(
    r'^\[\s*(\d+)\]\s+thr=0x([0-9a-fA-F]+)\s+(\S+)\s+'
    r'addr1=0x([0-9a-fA-F]+)\s+addr2=0x([0-9a-fA-F]+)\s+'
    r'data=(\d+)'
)
\end{lstlisting}

\noindent
Each match yields a dictionary with keys \texttt{timestamp},
\texttt{thread\_id}, \texttt{type}, \texttt{addr1}, \texttt{addr2},
and \texttt{data}.  Types include \texttt{TX\_BEGIN},
\texttt{COMMIT\_SUCCESS}, \texttt{TX\_ABORT}, \texttt{READ},
\texttt{WRITE}, \texttt{LOCK\_ACQUIRE}, and \texttt{LOCK\_RELEASE}.
A companion function \texttt{window\_around\_failure()} truncates the
log to the last $N$ events before an invariant violation, keeping
analysis tractable for logs exceeding 100\,000 lines.

\paragraph{Invariant checking.}
The file \texttt{invariant\_checker.py}~(\rep{tools/stm\_bug\_tool/invariant\_checker.py})
applies a suite of structural checks to the parsed events.  Each check
returns a list of violations:

\begin{lstlisting}[language=Python, caption={Lifecycle checker signature.}]
def check_complete_tx_lifecycle(parsed: dict) -> list:
    """Every TX_BEGIN must end with COMMIT_SUCCESS or TX_ABORT."""
\end{lstlisting}

\noindent
Other checks include lock discipline (every acquire must have a matching
release), read-write consistency (reads return values from the most
recent committed write to the same address), and absence of duplicate
reads within a single transaction.  The typical workflow is:

\begin{enumerate}
    \item Compile the benchmark with \texttt{-DTM\_EVENT\_LOGGER -O0}.
    \item Run the binary; on invariant failure it dumps the ring buffer to
          \texttt{stderr}.
    \item Pipe the output through the bug tool: \texttt{./my\_benchmark
          2\&1 | python3 tools/stm\_bug\_tool/invariant\_checker.py}.
    \item The tool prints violations with event ranges, narrowing the
          search to a handful of operations.
\end{enumerate}

The tool is backend-agnostic: the event logger is compiled into the
common runtime, so the same commands work for NOrec, TinySTM, TL2, or
any other backend.

\subsection{Debug Patches}
\index{debug patches}\index{patches/debug/}
\label{sec:debug-patches}

Debug output is useful during development but harmful in production:
\texttt{fprintf} calls on the hot path slow down benchmarks and clutter
CI logs.  The companion's solution is to keep debug output in
\emph{patches}, not in source.  The directory
\texttt{patches/debug/patches/}~(\rep{patches/debug/patches/})
contains numbered patch files that each add one category of diagnostics:

\begin{table}[htbp]
\centering
\caption{Debug patches and their purpose.}
\label{tab:debug-patches}
\begin{tabular}{M{0.32\linewidth}M{0.55\linewidth}}
\hline
\thb{0.32\textwidth}{Patch} & \thb{0.55\textwidth}{What it adds} \\
\hline
\texttt{002-tinystm-rs-ws-debug.patch} & Max read-set and write-set sizes on transaction exit \\
\texttt{003-tinystm-wbctl-debug ASSERT} & Lock-version mismatch diagnostics \\
\texttt{004-tm-region-alloc-debug.patch} & Mmap region base and size on allocation \\
\texttt{005-test-stress-ds-debug.patch} & Red-black tree pointer tracing \\
\texttt{006-swisstm-stats.patch} & SwissTM begin/end counts via \texttt{atexit} \\
\hline
\end{tabular}
\end{table}

\noindent
Applying all patches at once:

\begin{lstlisting}[language=bash, caption={Applying debug patches.}]
patches/debug/apply.sh     # applies all *.patch files
# ... run the backend under test ...
git checkout -- .           # or: patches/debug/revert.sh
\end{lstlisting}

\noindent
The \texttt{apply.sh} script refuses to run over a dirty working tree,
preventing accidental interleaving with uncommitted changes.  The
principle is simple: \emph{debug output should live in patches, not in
source}.  A developer debugging a TinySTM regression applies patch
002, rebuilds, reproduces, and reverts---the production source is
never touched.

\subsection{TSX Profiling}
\index{TSX profiling}\index{RDTSC instrumentation}
\label{sec:tsx-profiling}

Understanding performance requires hardware-level measurement.  The
companion's TSX profiling infrastructure instruments the TSXSGL
backend with RDTSC timestamps at each TSX primitive and produces
calibration data for the simulator's cost model.

The profiling patch (\rep{patches/profile/tsx/0001-tsxsgl-tsx-timing-instrumentation.patch})
adds cycle-level counters for: successful and failed \texttt{\_xbegin()}
calls, \texttt{\_xend()} commit time, explicit \texttt{\_xabort()}
penalty, SGL fallback lock/unlock/spin cycles, and per-operation read
and write latencies.  The workflow is fully scripted:

\begin{lstlisting}[language=bash, caption={The profiling workflow.}]
cd patches/profile/tsx
bash run_workflow.sh tsxsgl   # 1. apply patch
                       # 2. build
                       # 3. run experiments
                       # 4. revert patch
\end{lstlisting}

\noindent
The experiment runner \texttt{run\_tsx\_profiling.py}~(\rep{patches/profile/tsx/run\_tsx\_profiling.py})
executes \texttt{fuzz\_counter} and the bank benchmark across thread
counts 1, 2, 4, and 8 under both high-contention (64 accounts) and
low-contention (512 accounts) configurations.  It parses the
\texttt{TSX\_STATS} output lines, writes per-benchmark CSV results to
\texttt{profiling\_data/raw/}, and produces a calibration JSON that
maps each TSX event type to a measured cycle cost.

The calibration JSON feeds directly into the simulator's cost model.
The key insight is that the
profiling data captures real hardware costs on the target machine; the
simulator can then replay event traces with cycle-accurate timing
without re-running the instrumented binary.  The resulting
\texttt{machine\_profile.json} is portable: collect it on one machine,
simulate on another.

\subsection{LLVM Race Checker}
\index{race checker}\index{TM annotation check}
\label{sec:race-checker}

The instrumentation pass (\S\ref{sec:instrumentation-contract}) inserts
\texttt{tm\_read\_*/tm\_write\_*} calls for every access to a
TM-annotated global.  But non-transactional helper functions that
access the same globals are never instrumented---they execute outside
the transaction boundary and silently bypass conflict detection.

The race checker is a standalone LLVM pass
(\rep{plugin/passes/TMRaceCheckerPass.cpp}) that scans all
non-transactional functions for loads and stores to TM-annotated
globals and emits source-location warnings suggesting
\texttt{[[tm::shared]]} annotation:

\begin{lstlisting}[language=bash, caption={Running the race checker.}]
opt-22 -load-pass-plugin=plugin/bin/libTMRaceChecker.so \
       -passes="tm-race-checker" myapp.bc -o /dev/null
\end{lstlisting}

\noindent
The pass reuses \texttt{tracesFromTMGlobal()} from the instrumentation
pipeline---the same analysis that determines which variables need
instrumentation---and applies it in reverse: instead of instrumenting
accesses, it reports them.  A typical session is:

\begin{enumerate}
    \item Compile the application to LLVM IR: \texttt{clang-tm
          -emit-llvm -c myapp.cpp -o myapp.bc}.
    \item Run the race checker on the IR.
    \item For each warning, decide whether the access should be inside a
          transaction (annotate the caller with
          \texttt{[[tm::transaction\_al]]}) or whether the variable
          should be shared explicitly (annotate with
          \texttt{[[tm::shared]]}).
    \item Rebuild with the annotations and re-run the race checker to
          confirm the warning is resolved.
\end{enumerate}

The checker runs in under a second on typical benchmarks and integrates
into the build via the instrumentation plugin, which also calls
\texttt{checkMissingTransactionAnnotations()} at build time to warn
about missing annotations without a separate \texttt{opt} invocation.

\paragraph{Putting the tools together.}
A typical debugging session combines all four tools.  A bank test fails
with ``Money destroyed by 32'' on the NOrec backend.  The event logger
pinpoints a two-address write-back where one thread's commit overwrote
another's.  Debug patches reveal the write-set size was at the
capacity threshold, causing a partial rollback.  TSX profiling
(against the TL2 backend on the same workload) shows that the
conflict-avoidance benefit of back-off outweighs the validation cost.
The race checker confirms that all TM-annotated globals are properly
instrumented.  The failure is understood, the fix is targeted, and no
production source was modified along the way.
